% Source problem: https://www.ap.t.u-tokyo.ac.jp/themes/ap_t_u-tokyo/graduate/pdf/h24ap_exam.pdf
\documentclass[11pt,a4paper]{article}
\usepackage[margin=24mm]{geometry}
\usepackage{fontspec}
\usepackage{xeCJK}
\setCJKmainfont{HaranoAjiMincho-Regular.otf}
\setCJKsansfont{HaranoAjiGothic-Medium.otf}
\usepackage{amsmath,amssymb}
\usepackage{tikz}
\usetikzlibrary{arrows.meta}
\setlength{\parindent}{1em}
\setlength{\parskip}{0.35em}
\renewcommand{\figurename}{図}
\title{物理工学専攻 院試解答集\\2012年度 物理学II 第3問}
\author{政岡凜太郎}
\date{}
\begin{document}
\maketitle

全設問で時間依存を $e^{-i\omega t}$ とする。$\boldsymbol H$ は $y$ 成分のみを持ち、電場は $x,z$ 成分を持つTM波である。媒質 $j$ では $\boldsymbol D=\epsilon_j\boldsymbol E$、$\boldsymbol B=\mu_0\boldsymbol H$。

\section*{[1] 全反射と真空中の減衰波}
電場は電荷のない一様媒質中で $\operatorname{div}\boldsymbol E=0$ を満たす。Faraday則 $\operatorname{rot}\boldsymbol E=i\omega\boldsymbol B$ と Amp\`ere則 $\operatorname{rot}\boldsymbol H=-i\omega\boldsymbol D$ を用いると、
\begin{align}
  \operatorname{rot}(\operatorname{rot}\boldsymbol E)
  &=i\omega\mu_0\operatorname{rot}\boldsymbol H
   =\mu_0\epsilon_j\omega^2\boldsymbol E,\\
  -\Delta\boldsymbol E&=\mu_0\epsilon_j\omega^2\boldsymbol E.
\end{align}
よって平面波の波数は $k_x^2+k_z^2=\epsilon_j\mu_0\omega^2$ に従う。媒質1では $\epsilon_1\mu_0=n^2/c^2$ だから、入射波の $z$ 成分は
\begin{equation}
  \boxed{k_1=\sqrt{\frac{n^2\omega^2}{c^2}-k_x^2}}.
\end{equation}
真空側の場 $e^{ik_xx-\kappa_2z}$ を平面波の式に代入すると、$z$ 波数は $i\kappa_2$ であるため
\begin{equation}
  k_x^2-\kappa_2^2=\frac{\omega^2}{c^2},\qquad
  \boxed{\kappa_2=\sqrt{k_x^2-\frac{\omega^2}{c^2}}}.
\end{equation}
入射側に伝播波があり $k_1>0$、透過側で減衰するため $\kappa_2>0$ でなければならない。従って
\begin{equation}
  \boxed{\frac{\omega}{c}<k_x<\frac{n\omega}{c}}.
\end{equation}
下端は真空側の減衰条件、上端は媒質1側の入射波の伝播条件から来る。端点ではそれぞれ $\kappa_2=0$ または $k_1=0$ となるので、指定された狭義の条件には含まれない。

\section*{[2] 媒質間の無反射透過}
平面波 $e^{i(k_xx+k_jz)}$ では $\boldsymbol k_j=(k_x,0,k_j)$ である。$\operatorname{div}\boldsymbol D=0$ より
\begin{equation}
  k_xE_{j,x}+k_jE_{j,z}=0,
  \qquad E_{j,z}=-\frac{k_x}{k_j}E_{j,x}.
\end{equation}
また $\operatorname{rot}\boldsymbol E=i\omega\mu_0\boldsymbol H$ の $y$ 成分をとると
\begin{align}
  k_jE_{j,x}-k_xE_{j,z}&=\omega\mu_0 H_{j,y},\\
  \omega\mu_0 H_{j,y}&=\frac{k_x^2+k_j^2}{k_j}E_{j,x}
   =\frac{\epsilon_j\mu_0\omega^2}{k_j}E_{j,x}.
\end{align}
したがって各媒質の電場・磁場比は
\begin{equation}
  \boxed{\zeta_j=\frac{E_{j,x}}{H_{j,y}}=\frac{k_j}{\epsilon_j\omega}}.
\end{equation}
反射がないとき接線成分 $E_x,H_y$ は界面で連続なので $\zeta_1=\zeta_2$、すなわち $k_1/\epsilon_1=k_2/\epsilon_2$ である。分散関係
\begin{equation}
  k_1^2=\epsilon_1\mu_0\omega^2-k_x^2,\qquad
  k_2^2=\epsilon_2\mu_0\omega^2-k_x^2
\end{equation}
を用いると
\begin{align}
  \frac{\epsilon_1\mu_0\omega^2-k_x^2}{\epsilon_1^2}
  &=\frac{\epsilon_2\mu_0\omega^2-k_x^2}{\epsilon_2^2},\\
  \boxed{k_x^2=\frac{\mu_0\omega^2\epsilon_1\epsilon_2}{\epsilon_1+\epsilon_2}},\qquad
  \boxed{k_x=\omega\sqrt{\frac{\mu_0\epsilon_1\epsilon_2}{\epsilon_1+\epsilon_2}}}.
\end{align}
この波数が入射側の伝播条件 $k_x^2<\epsilon_1\mu_0\omega^2$ と透過側の条件 $k_x^2<\epsilon_2\mu_0\omega^2$ を満たすことを確認できる。実際、$\epsilon_2>\epsilon_1>0$ のとき比はそれぞれ $\epsilon_2/(\epsilon_1+\epsilon_2)<1$、$\epsilon_1/(\epsilon_1+\epsilon_2)<1$ であり、$k_1,k_2$ は正の実数となる。接線電場・磁場を同時に連続にするこの入射成分では、反射波の振幅はゼロとなる。

\section*{[3] 金属界面に束縛された波}
$+x$ へ伝わる枝を選び、$k_x>0$ とする。逆向きの枝は同じ分散を持つ $-k_x$ で表される。
$z<0$ の減衰因子は $e^{+\kappa_1z}=e^{ik_{z,1}z}$ に対応し、$k_{z,1}=-i\kappa_1$。一方、$z>0$ の因子 $e^{-\kappa_2z}$ は $k_{z,2}=+i\kappa_2$ に対応する。各媒質の分散関係から
\begin{equation}
  \kappa_1^2=k_x^2-\epsilon_0\mu_0\omega^2,\qquad
  \kappa_2^2=k_x^2-\epsilon_2\mu_0\omega^2
  =k_x^2+|\epsilon_2|\mu_0\omega^2.
\end{equation}
[2]で導いた $E_x/H_y=k_z/(\epsilon\omega)$ を両側に適用し、界面で接線成分を連続にすると
\begin{equation}
  -\frac{i\kappa_1}{\epsilon_0\omega}
  =\frac{i\kappa_2}{\epsilon_2\omega},\qquad
  \boxed{\frac{\kappa_1}{\kappa_2}=\frac{\epsilon_0}{|\epsilon_2|}}.
\end{equation}
この関係と分散関係を組み合わせれば
\begin{align}
  \frac{k_x^2-\epsilon_0\mu_0\omega^2}
       {k_x^2+|\epsilon_2|\mu_0\omega^2}
  &=\frac{\epsilon_0^2}{|\epsilon_2|^2},\\
  \boxed{k_x^2} &=\frac{\mu_0\omega^2\epsilon_0|\epsilon_2|}
                        {|\epsilon_2|-\epsilon_0},\\
  \boxed{k_x} &=\omega\sqrt{\frac{\mu_0\epsilon_0|\epsilon_2|}
                                      {|\epsilon_2|-\epsilon_0}}.
\end{align}
$|\epsilon_2|>\epsilon_0$ なので
\begin{equation}
  \boxed{k_x=\frac{\omega}{c}\sqrt{\frac{|\epsilon_2|}{|\epsilon_2|-\epsilon_0}}>
  \frac{\omega}{c}}.
\end{equation}
この面内波数は真空中の伝播光の最大値 $\omega/c$ より大きい。したがってこのモードは界面に沿って進む一方、垂直方向には真空側・金属側の双方へ減衰する束縛波である。

\section*{[4] プリズムによる励起}
真空中から金属表面へ直接入射する伝播光は、界面に平行な波数が $k_x\leq\omega/c$ である。一方、[3]の表面波は $k_x>\omega/c$ を要するので、界面に平行な運動量が足りず、平坦な面では直接結合できない。

屈折率 $n_{\mathrm p}>1$ のプリズムから界面の法線となす入射角 $\theta$ で光を入れると、界面に平行な波数は $k_x=n_{\mathrm p}(\omega/c)\sin\theta$ まで大きくなる。真空層との境界で全反射する角度を選べば、この高い面内波数を保ったエバネッセント場が真空ギャップへしみ出す。ギャップを金属表面の近くに狭く保つことで場が金属へ届き、さらに $k_x$ を[3]の表面波波数に合わせればそのモードを励起できる。つまりプリズムは、全反射場を介して伝播光では得られない面内波数を供給し、金属界面の束縛モードとの運動量整合を可能にする。

\begin{figure}[htbp]
\centering
\begin{tikzpicture}[>=Stealth,scale=1.0]
  \pgfmathsetmacro{\incidenceAngle}{55}
  \pgfmathsetmacro{\rayLength}{1.05}
  \pgfmathsetmacro{\hitX}{3.3}
  \pgfmathsetmacro{\hitZ}{0.3}
  \pgfmathsetmacro{\gapWidth}{0.5}
  \pgfmathsetmacro{\metalX}{\hitX+\gapWidth}
  \pgfmathsetmacro{\surfaceArrowX}{\metalX-0.04}
  \pgfmathsetmacro{\rayStartX}{\hitX-\rayLength*cos(\incidenceAngle)}
  \pgfmathsetmacro{\rayStartZ}{\hitZ-\rayLength*sin(\incidenceAngle)}
  \pgfmathsetmacro{\rayEndX}{\hitX-\rayLength*cos(\incidenceAngle)}
  \pgfmathsetmacro{\rayEndZ}{\hitZ+\rayLength*sin(\incidenceAngle)}
  \draw[thick,fill=gray!15] (0,0) -- (3.3,1.5) -- (3.3,-1.5) -- cycle;
  \draw[very thick] (\metalX,-1.6) -- (\metalX,1.6);
  \node at (1.5,0.2) {プリズム};
  \node[anchor=west] at (\metalX+0.25,1.25) {金属};
  \node at (\hitX+\gapWidth/2,1.8) {真空ギャップ};
  \draw[->,thick] (\rayStartX,\rayStartZ) -- (\hitX,\hitZ) node[midway,left] {入射光};
  \draw[->,thick] (\hitX,\hitZ) -- (\rayEndX,\rayEndZ);
  \draw[thick,dashed] (\hitX,\hitZ) -- (\metalX,\hitZ);
  \draw[->,thick,blue] (\surfaceArrowX,-0.45) -- (\surfaceArrowX,1.0);
  \node[anchor=west] at (\metalX+0.3,0.25) {界面に沿う表面波};
  \draw[<->] (\hitX,-1.85) -- (\metalX,-1.85) node[midway,below] {狭い隙間};
\end{tikzpicture}
\caption{全反射によるエバネッセント場を介したプリズム結合の模式図。点線はギャップ内にしみ出す場、青い矢印は金属界面に平行な表面波の進行方向を表す。入射角と長さは説明用で、臨界角や寸法を表さない。}
\end{figure}

\end{document}
